% ####################################################################
% Basic configuration of Beamera class and Jagiellonian theme
% ####################################################################
\RequirePackage[l2tabu, orthodox]{nag}



\ifx\PresentationStyle\notset
  \def\PresentationStyle{dark}
\fi



% Options: t - align frame text to the top.
\documentclass[10pt,t]{beamer}
\mode<presentation>
\usetheme[style=\PresentationStyle,JUlogotitle=no]{jagiellonian}




% ------------------------------------------------------------------------------------
% Procesing configuration files of Jagiellonian theme located
% in the directory "preambule"
% ------------------------------------------------------------------------------------
% Configuration for polish language
% Need description
\usepackage[polish]{babel}
% Need description
\usepackage[MeX]{polski}



% ------------------------------
% Better support of polish chars in technical parts of PDF
% ------------------------------
\hypersetup{pdfencoding=auto,psdextra}

% Package "textpos" give as enviroment "textblock" which is very usefull in
% arranging text on slides.

% This is standard configuration of "textpos"
\usepackage[overlay,absolute]{textpos}

% If you need to see bounds of "textblock's" comment line above and uncomment
% one below.

% Caution! When showboxes option is on significant ammunt of space is add
% to the top of textblock and as such, everyting put in them gone down.
% We need to check how to remove this bug.

% \usepackage[showboxes,overlay,absolute]{textpos}



% Setting scale length for package "textpos"
\setlength{\TPHorizModule}{10mm}
\setlength{\TPVertModule}{\TPHorizModule}


% ---------------------------------------
% Packages written for lectures "Geometria 3D dla twórców gier wideo"
% ---------------------------------------
% \usepackage{./Geometry3DPackages/Geometry3D}
% \usepackage{./Geometry3DPackages/Geometry3DIndices}
% \usepackage{./Geometry3DPackages/Geometry3DTikZStyle}
% \usepackage{./ProgramowanieSymulacjiFizykiPaczki/ProgramowanieSymulacjiFizykiTikZStyle}
% \usepackage{./Geometry3DPackages/mathcommands}


% ---------------------------------------
% Importing TikZ's libraries
% ---------------------------------------
\usetikzlibrary{arrows.meta}
\usetikzlibrary{positioning}





% % Configuration package "bm" that need for making bold symbols
\newcommand{\bmmax}{0}
\newcommand{\hmmax}{0}
\usepackage{bm}




% ---------------------------------------
% Packages for scientific texts
% ---------------------------------------
% \let\lll\undefined  % Sometimes you must use this line to allow
% "amsmath" package to works with packages with packages for polish
% languge imported
% /preambul/LanguageSettings/JagiellonianPolishLanguageSettings.tex.
% This comments (probably) removes polish letter Ł.
\usepackage{amsmath}  % Packages from American Mathematical Society (AMS)
\usepackage{amssymb}
\usepackage{amscd}
\usepackage{amsthm}
\usepackage{siunitx}  % Package for typsetting SI units.
\usepackage{upgreek}  % Better looking greek letters.
% Example of using upgreek: pi = \uppi


\usepackage{calrsfs}  % Zmienia czcionkę kaligraficzną w \mathcal
% na ładniejszą. Może w innych miejscach robi to samo, ale o tym nic
% nie wiem.







% !!!!!!!!!!!!!!!!!!!!!!!!!!!!!!
% !!!!!!!!!!!!!!!!!!!!!!!!!!!!!!
% EVIL STUFF
\if\JUlogotitle1
\edef\LogoJUPath{LogoJU_\JUlogoLang/LogoJU_\JUlogoShape_\JUlogoColor.pdf}
\titlegraphic{\hfill\includegraphics[scale=0.22]
{./JagiellonianPictures/\LogoJUPath}}
\fi

% ---------------------------------------
% Commands for handling colors
% ---------------------------------------


% Command for setting normal text color for some text in math modestyle
% Text color depend on used style of Jagiellonian

% Beamer version of command
\newcommand{\TextWithNormalTextColor}[1]{%
  {\color{jNormalTextFGColor}
    \setbeamercolor{math text}{fg=jNormalTextFGColor} {#1}}
}

% Article and similar classes version of command
% \newcommand{\TextWithNormalTextColor}[1]{%
%   {\color{jNormalTextsFGColor} {#1}}
% }



% Beamer version of command
\newcommand{\NormalTextInMathMode}[1]{%
  {\color{jNormalTextFGColor}
    \setbeamercolor{math text}{fg=jNormalTextFGColor} \text{#1}}
}


% Article and similar classes version of command
% \newcommand{\NormalTextInMathMode}[1]{%
%   {\color{jNormalTextsFGColor} \text{#1}}
% }




% Command that sets color of some mathematical text to the same color
% that has normal text in header (?)

% Beamer version of the command
\newcommand{\MathTextFrametitleFGColor}[1]{%
  {\color{jFrametitleFGColor}
    \setbeamercolor{math text}{fg=jFrametitleFGColor} #1}
}

% Article and similar classes version of the command
% \newcommand{\MathTextWhiteColor}[1]{{\color{jFrametitleFGColor} #1}}





% Command for setting color of alert text for some text in math modestyle

% Beamer version of the command
\newcommand{\MathTextAlertColor}[1]{%
  {\color{jOrange} \setbeamercolor{math text}{fg=jOrange} #1}
}

% Article and similar classes version of the command
% \newcommand{\MathTextAlertColor}[1]{{\color{jOrange} #1}}





% Command that allow you to sets chosen color as the color of some text into
% math mode. Due to some nuances in the way that Beamer handle colors
% it not work in all cases. We hope that in the future we will improve it.

% Beamer version of the command
\newcommand{\SetMathTextColor}[2]{%
  {\color{#1} \setbeamercolor{math text}{fg=#1} #2}
}


% Article and similar classes version of the command
% \newcommand{\SetMathTextColor}[2]{{\color{#1} #2}}










% ---------------------------------------
% Commands for few special slides
% ---------------------------------------
\newcommand{\EndingSlide}[1]{%
  \begin{frame}[standout]

    \begingroup

    \color{jFrametitleFGColor}

    #1

    \endgroup

  \end{frame}
}










% ---------------------------------------
% Commands for setting background pictures for some slides
% ---------------------------------------
\newcommand{\TitleBackgroundPicture}
{./JagiellonianPictures/Backgrounds/LajkonikDark.png}
\newcommand{\SectionBackgroundPicture}
{./JagiellonianPictures/Backgrounds/LajkonikLight.png}



\newcommand{\TitleSlideWithPicture}{%
  \begingroup

  \usebackgroundtemplate{%
    \includegraphics[height=\paperheight]{\TitleBackgroundPicture}}

  \maketitle

  \endgroup
}





\newcommand{\SectionSlideWithPicture}[1]{%
  \begingroup

  \usebackgroundtemplate{%
    \includegraphics[height=\paperheight]{\SectionBackgroundPicture}}

  \setbeamercolor{titlelike}{fg=normal text.fg}

  \section{#1}

  \endgroup
}










% ---------------------------------------
% Commands for lectures "Geometria 3D dla twórców gier wideo"
% Polish version
% ---------------------------------------
% Komendy teraz wykomentowane były potrzebne, gdy loga były na niebieskim
% tle, nie na białym. A są na białym bo tego chcieli w biurze projektu.
% \newcommand{\FundingLogoWhitePicturePL}
% {./PresentationPictures/CommonPictures/logotypFundusze_biale_bez_tla2.pdf}
\newcommand{\FundingLogoColorPicturePL}
{./PresentationPictures/CommonPictures/European_Funds_color_PL.pdf}
% \newcommand{\EULogoWhitePicturePL}
% {./PresentationPictures/CommonPictures/logotypUE_biale_bez_tla2.pdf}
\newcommand{\EUSocialFundLogoColorPicturePL}
{./PresentationPictures/CommonPictures/EU_Social_Fund_color_PL.pdf}
% \newcommand{\ZintegrUJLogoWhitePicturePL}
% {./PresentationPictures/CommonPictures/zintegruj-logo-white.pdf}
\newcommand{\ZintegrUJLogoColorPicturePL}
{./PresentationPictures/CommonPictures/ZintegrUJ_color.pdf}
\newcommand{\JULogoColorPicturePL}
{./JagiellonianPictures/LogoJU_PL/LogoJU_A_color.pdf}





\newcommand{\GeometryThreeDSpecialBeginningSlidePL}{%
  \begin{frame}[standout]

    \begin{textblock}{11}(1,0.7)

      \begin{flushleft}

        \mdseries

        \footnotesize

        \color{jFrametitleFGColor}

        Materiał powstał w ramach projektu współfinansowanego ze środków
        Unii Europejskiej w ramach Europejskiego Funduszu Społecznego
        POWR.03.05.00-00-Z309/17-00.

      \end{flushleft}

    \end{textblock}





    \begin{textblock}{10}(0,2.2)

      \tikz \fill[color=jBackgroundStyleLight] (0,0) rectangle (12.8,-1.5);

    \end{textblock}


    \begin{textblock}{3.2}(1,2.45)

      \includegraphics[scale=0.3]{\FundingLogoColorPicturePL}

    \end{textblock}


    \begin{textblock}{2.5}(3.7,2.5)

      \includegraphics[scale=0.2]{\JULogoColorPicturePL}

    \end{textblock}


    \begin{textblock}{2.5}(6,2.4)

      \includegraphics[scale=0.1]{\ZintegrUJLogoColorPicturePL}

    \end{textblock}


    \begin{textblock}{4.2}(8.4,2.6)

      \includegraphics[scale=0.3]{\EUSocialFundLogoColorPicturePL}

    \end{textblock}

  \end{frame}
}



\newcommand{\GeometryThreeDTwoSpecialBeginningSlidesPL}{%
  \begin{frame}[standout]

    \begin{textblock}{11}(1,0.7)

      \begin{flushleft}

        \mdseries

        \footnotesize

        \color{jFrametitleFGColor}

        Materiał powstał w ramach projektu współfinansowanego ze środków
        Unii Europejskiej w ramach Europejskiego Funduszu Społecznego
        POWR.03.05.00-00-Z309/17-00.

      \end{flushleft}

    \end{textblock}





    \begin{textblock}{10}(0,2.2)

      \tikz \fill[color=jBackgroundStyleLight] (0,0) rectangle (12.8,-1.5);

    \end{textblock}


    \begin{textblock}{3.2}(1,2.45)

      \includegraphics[scale=0.3]{\FundingLogoColorPicturePL}

    \end{textblock}


    \begin{textblock}{2.5}(3.7,2.5)

      \includegraphics[scale=0.2]{\JULogoColorPicturePL}

    \end{textblock}


    \begin{textblock}{2.5}(6,2.4)

      \includegraphics[scale=0.1]{\ZintegrUJLogoColorPicturePL}

    \end{textblock}


    \begin{textblock}{4.2}(8.4,2.6)

      \includegraphics[scale=0.3]{\EUSocialFundLogoColorPicturePL}

    \end{textblock}

  \end{frame}





  \TitleSlideWithPicture
}



\newcommand{\GeometryThreeDSpecialEndingSlidePL}{%
  \begin{frame}[standout]

    \begin{textblock}{11}(1,0.7)

      \begin{flushleft}

        \mdseries

        \footnotesize

        \color{jFrametitleFGColor}

        Materiał powstał w ramach projektu współfinansowanego ze środków
        Unii Europejskiej w~ramach Europejskiego Funduszu Społecznego
        POWR.03.05.00-00-Z309/17-00.

      \end{flushleft}

    \end{textblock}





    \begin{textblock}{10}(0,2.2)

      \tikz \fill[color=jBackgroundStyleLight] (0,0) rectangle (12.8,-1.5);

    \end{textblock}


    \begin{textblock}{3.2}(1,2.45)

      \includegraphics[scale=0.3]{\FundingLogoColorPicturePL}

    \end{textblock}


    \begin{textblock}{2.5}(3.7,2.5)

      \includegraphics[scale=0.2]{\JULogoColorPicturePL}

    \end{textblock}


    \begin{textblock}{2.5}(6,2.4)

      \includegraphics[scale=0.1]{\ZintegrUJLogoColorPicturePL}

    \end{textblock}


    \begin{textblock}{4.2}(8.4,2.6)

      \includegraphics[scale=0.3]{\EUSocialFundLogoColorPicturePL}

    \end{textblock}





    \begin{textblock}{11}(1,4)

      \begin{flushleft}

        \mdseries

        \footnotesize

        \RaggedRight

        \color{jFrametitleFGColor}

        Treść niniejszego wykładu jest udostępniona na~licencji
        Creative Commons (\textsc{cc}), z~uzna\-niem autorstwa
        (\textsc{by}) oraz udostępnianiem na tych samych warunkach
        (\textsc{sa}). Rysunki i~wy\-kresy zawarte w~wykładzie są
        autorstwa dr.~hab.~Pawła Węgrzyna et~al. i~są dostępne
        na tej samej licencji, o~ile nie wskazano inaczej.
        W~prezentacji wykorzystano temat Beamera Jagiellonian,
        oparty na~temacie Metropolis Matthiasa Vogelgesanga,
        dostępnym na licencji \LaTeX{} Project Public License~1.3c
        pod adresem: \colorhref{https://github.com/matze/mtheme}
        {https://github.com/matze/mtheme}.

        Projekt typograficzny: Iwona Grabska-Gradzińska \\
        Skład: Kamil Ziemian;
        Korekta: Wojciech Palacz \\
        Modele: Dariusz Frymus, Kamil Nowakowski \\
        Rysunki i~wykresy: Kamil Ziemian, Paweł Węgrzyn, Wojciech Palacz

      \end{flushleft}

    \end{textblock}

  \end{frame}
}



\newcommand{\GeometryThreeDTwoSpecialEndingSlidesPL}[1]{%
  \begin{frame}[standout]


    \begin{textblock}{11}(1,0.7)

      \begin{flushleft}

        \mdseries

        \footnotesize

        \color{jFrametitleFGColor}

        Materiał powstał w ramach projektu współfinansowanego ze środków
        Unii Europejskiej w~ramach Europejskiego Funduszu Społecznego
        POWR.03.05.00-00-Z309/17-00.

      \end{flushleft}

    \end{textblock}





    \begin{textblock}{10}(0,2.2)

      \tikz \fill[color=jBackgroundStyleLight] (0,0) rectangle (12.8,-1.5);

    \end{textblock}


    \begin{textblock}{3.2}(1,2.45)

      \includegraphics[scale=0.3]{\FundingLogoColorPicturePL}

    \end{textblock}


    \begin{textblock}{2.5}(3.7,2.5)

      \includegraphics[scale=0.2]{\JULogoColorPicturePL}

    \end{textblock}


    \begin{textblock}{2.5}(6,2.4)

      \includegraphics[scale=0.1]{\ZintegrUJLogoColorPicturePL}

    \end{textblock}


    \begin{textblock}{4.2}(8.4,2.6)

      \includegraphics[scale=0.3]{\EUSocialFundLogoColorPicturePL}

    \end{textblock}





    \begin{textblock}{11}(1,4)

      \begin{flushleft}

        \mdseries

        \footnotesize

        \RaggedRight

        \color{jFrametitleFGColor}

        Treść niniejszego wykładu jest udostępniona na~licencji
        Creative Commons (\textsc{cc}), z~uzna\-niem autorstwa
        (\textsc{by}) oraz udostępnianiem na tych samych warunkach
        (\textsc{sa}). Rysunki i~wy\-kresy zawarte w~wykładzie są
        autorstwa dr.~hab.~Pawła Węgrzyna et~al. i~są dostępne
        na tej samej licencji, o~ile nie wskazano inaczej.
        W~prezentacji wykorzystano temat Beamera Jagiellonian,
        oparty na~temacie Metropolis Matthiasa Vogelgesanga,
        dostępnym na licencji \LaTeX{} Project Public License~1.3c
        pod adresem: \colorhref{https://github.com/matze/mtheme}
        {https://github.com/matze/mtheme}.

        Projekt typograficzny: Iwona Grabska-Gradzińska \\
        Skład: Kamil Ziemian;
        Korekta: Wojciech Palacz \\
        Modele: Dariusz Frymus, Kamil Nowakowski \\
        Rysunki i~wykresy: Kamil Ziemian, Paweł Węgrzyn, Wojciech Palacz

      \end{flushleft}

    \end{textblock}

  \end{frame}





  \begin{frame}[standout]

    \begingroup

    \color{jFrametitleFGColor}

    #1

    \endgroup

  \end{frame}
}



\newcommand{\GeometryThreeDSpecialEndingSlideVideoPL}{%
  \begin{frame}[standout]

    \begin{textblock}{11}(1,0.7)

      \begin{flushleft}

        \mdseries

        \footnotesize

        \color{jFrametitleFGColor}

        Materiał powstał w ramach projektu współfinansowanego ze środków
        Unii Europejskiej w~ramach Europejskiego Funduszu Społecznego
        POWR.03.05.00-00-Z309/17-00.

      \end{flushleft}

    \end{textblock}





    \begin{textblock}{10}(0,2.2)

      \tikz \fill[color=jBackgroundStyleLight] (0,0) rectangle (12.8,-1.5);

    \end{textblock}


    \begin{textblock}{3.2}(1,2.45)

      \includegraphics[scale=0.3]{\FundingLogoColorPicturePL}

    \end{textblock}


    \begin{textblock}{2.5}(3.7,2.5)

      \includegraphics[scale=0.2]{\JULogoColorPicturePL}

    \end{textblock}


    \begin{textblock}{2.5}(6,2.4)

      \includegraphics[scale=0.1]{\ZintegrUJLogoColorPicturePL}

    \end{textblock}


    \begin{textblock}{4.2}(8.4,2.6)

      \includegraphics[scale=0.3]{\EUSocialFundLogoColorPicturePL}

    \end{textblock}





    \begin{textblock}{11}(1,4)

      \begin{flushleft}

        \mdseries

        \footnotesize

        \RaggedRight

        \color{jFrametitleFGColor}

        Treść niniejszego wykładu jest udostępniona na~licencji
        Creative Commons (\textsc{cc}), z~uzna\-niem autorstwa
        (\textsc{by}) oraz udostępnianiem na tych samych warunkach
        (\textsc{sa}). Rysunki i~wy\-kresy zawarte w~wykładzie są
        autorstwa dr.~hab.~Pawła Węgrzyna et~al. i~są dostępne
        na tej samej licencji, o~ile nie wskazano inaczej.
        W~prezentacji wykorzystano temat Beamera Jagiellonian,
        oparty na~temacie Metropolis Matthiasa Vogelgesanga,
        dostępnym na licencji \LaTeX{} Project Public License~1.3c
        pod adresem: \colorhref{https://github.com/matze/mtheme}
        {https://github.com/matze/mtheme}.

        Projekt typograficzny: Iwona Grabska-Gradzińska;
        Skład: Kamil Ziemian \\
        Korekta: Wojciech Palacz;
        Modele: Dariusz Frymus, Kamil Nowakowski \\
        Rysunki i~wykresy: Kamil Ziemian, Paweł Węgrzyn, Wojciech Palacz \\
        Montaż: Agencja Filmowa Film \& Television Production~-- Zbigniew
        Masklak

      \end{flushleft}

    \end{textblock}

  \end{frame}
}





\newcommand{\GeometryThreeDTwoSpecialEndingSlidesVideoPL}[1]{%
  \begin{frame}[standout]

    \begin{textblock}{11}(1,0.7)

      \begin{flushleft}

        \mdseries

        \footnotesize

        \color{jFrametitleFGColor}

        Materiał powstał w ramach projektu współfinansowanego ze środków
        Unii Europejskiej w~ramach Europejskiego Funduszu Społecznego
        POWR.03.05.00-00-Z309/17-00.

      \end{flushleft}

    \end{textblock}





    \begin{textblock}{10}(0,2.2)

      \tikz \fill[color=jBackgroundStyleLight] (0,0) rectangle (12.8,-1.5);

    \end{textblock}


    \begin{textblock}{3.2}(1,2.45)

      \includegraphics[scale=0.3]{\FundingLogoColorPicturePL}

    \end{textblock}


    \begin{textblock}{2.5}(3.7,2.5)

      \includegraphics[scale=0.2]{\JULogoColorPicturePL}

    \end{textblock}


    \begin{textblock}{2.5}(6,2.4)

      \includegraphics[scale=0.1]{\ZintegrUJLogoColorPicturePL}

    \end{textblock}


    \begin{textblock}{4.2}(8.4,2.6)

      \includegraphics[scale=0.3]{\EUSocialFundLogoColorPicturePL}

    \end{textblock}





    \begin{textblock}{11}(1,4)

      \begin{flushleft}

        \mdseries

        \footnotesize

        \RaggedRight

        \color{jFrametitleFGColor}

        Treść niniejszego wykładu jest udostępniona na~licencji
        Creative Commons (\textsc{cc}), z~uzna\-niem autorstwa
        (\textsc{by}) oraz udostępnianiem na tych samych warunkach
        (\textsc{sa}). Rysunki i~wy\-kresy zawarte w~wykładzie są
        autorstwa dr.~hab.~Pawła Węgrzyna et~al. i~są dostępne
        na tej samej licencji, o~ile nie wskazano inaczej.
        W~prezentacji wykorzystano temat Beamera Jagiellonian,
        oparty na~temacie Metropolis Matthiasa Vogelgesanga,
        dostępnym na licencji \LaTeX{} Project Public License~1.3c
        pod adresem: \colorhref{https://github.com/matze/mtheme}
        {https://github.com/matze/mtheme}.

        Projekt typograficzny: Iwona Grabska-Gradzińska;
        Skład: Kamil Ziemian \\
        Korekta: Wojciech Palacz;
        Modele: Dariusz Frymus, Kamil Nowakowski \\
        Rysunki i~wykresy: Kamil Ziemian, Paweł Węgrzyn, Wojciech Palacz \\
        Montaż: Agencja Filmowa Film \& Television Production~-- Zbigniew
        Masklak

      \end{flushleft}

    \end{textblock}

  \end{frame}





  \begin{frame}[standout]


    \begingroup

    \color{jFrametitleFGColor}

    #1

    \endgroup

  \end{frame}
}










% ---------------------------------------
% Commands for lectures "Geometria 3D dla twórców gier wideo"
% English version
% ---------------------------------------
% \newcommand{\FundingLogoWhitePictureEN}
% {./PresentationPictures/CommonPictures/logotypFundusze_biale_bez_tla2.pdf}
\newcommand{\FundingLogoColorPictureEN}
{./PresentationPictures/CommonPictures/European_Funds_color_EN.pdf}
% \newcommand{\EULogoWhitePictureEN}
% {./PresentationPictures/CommonPictures/logotypUE_biale_bez_tla2.pdf}
\newcommand{\EUSocialFundLogoColorPictureEN}
{./PresentationPictures/CommonPictures/EU_Social_Fund_color_EN.pdf}
% \newcommand{\ZintegrUJLogoWhitePictureEN}
% {./PresentationPictures/CommonPictures/zintegruj-logo-white.pdf}
\newcommand{\ZintegrUJLogoColorPictureEN}
{./PresentationPictures/CommonPictures/ZintegrUJ_color.pdf}
\newcommand{\JULogoColorPictureEN}
{./JagiellonianPictures/LogoJU_EN/LogoJU_A_color.pdf}



\newcommand{\GeometryThreeDSpecialBeginningSlideEN}{%
  \begin{frame}[standout]

    \begin{textblock}{11}(1,0.7)

      \begin{flushleft}

        \mdseries

        \footnotesize

        \color{jFrametitleFGColor}

        This content was created as part of a project co-financed by the
        European Union within the framework of the European Social Fund
        POWR.03.05.00-00-Z309/17-00.

      \end{flushleft}

    \end{textblock}





    \begin{textblock}{10}(0,2.2)

      \tikz \fill[color=jBackgroundStyleLight] (0,0) rectangle (12.8,-1.5);

    \end{textblock}


    \begin{textblock}{3.2}(0.7,2.45)

      \includegraphics[scale=0.3]{\FundingLogoColorPictureEN}

    \end{textblock}


    \begin{textblock}{2.5}(4.15,2.5)

      \includegraphics[scale=0.2]{\JULogoColorPictureEN}

    \end{textblock}


    \begin{textblock}{2.5}(6.35,2.4)

      \includegraphics[scale=0.1]{\ZintegrUJLogoColorPictureEN}

    \end{textblock}


    \begin{textblock}{4.2}(8.4,2.6)

      \includegraphics[scale=0.3]{\EUSocialFundLogoColorPictureEN}

    \end{textblock}

  \end{frame}
}



\newcommand{\GeometryThreeDTwoSpecialBeginningSlidesEN}{%
  \begin{frame}[standout]

    \begin{textblock}{11}(1,0.7)

      \begin{flushleft}

        \mdseries

        \footnotesize

        \color{jFrametitleFGColor}

        This content was created as part of a project co-financed by the
        European Union within the framework of the European Social Fund
        POWR.03.05.00-00-Z309/17-00.

      \end{flushleft}

    \end{textblock}





    \begin{textblock}{10}(0,2.2)

      \tikz \fill[color=jBackgroundStyleLight] (0,0) rectangle (12.8,-1.5);

    \end{textblock}


    \begin{textblock}{3.2}(0.7,2.45)

      \includegraphics[scale=0.3]{\FundingLogoColorPictureEN}

    \end{textblock}


    \begin{textblock}{2.5}(4.15,2.5)

      \includegraphics[scale=0.2]{\JULogoColorPictureEN}

    \end{textblock}


    \begin{textblock}{2.5}(6.35,2.4)

      \includegraphics[scale=0.1]{\ZintegrUJLogoColorPictureEN}

    \end{textblock}


    \begin{textblock}{4.2}(8.4,2.6)

      \includegraphics[scale=0.3]{\EUSocialFundLogoColorPictureEN}

    \end{textblock}

  \end{frame}





  \TitleSlideWithPicture
}



\newcommand{\GeometryThreeDSpecialEndingSlideEN}{%
  \begin{frame}[standout]

    \begin{textblock}{11}(1,0.7)

      \begin{flushleft}

        \mdseries

        \footnotesize

        \color{jFrametitleFGColor}

        This content was created as part of a project co-financed by the
        European Union within the framework of the European Social Fund
        POWR.03.05.00-00-Z309/17-00.

      \end{flushleft}

    \end{textblock}





    \begin{textblock}{10}(0,2.2)

      \tikz \fill[color=jBackgroundStyleLight] (0,0) rectangle (12.8,-1.5);

    \end{textblock}


    \begin{textblock}{3.2}(0.7,2.45)

      \includegraphics[scale=0.3]{\FundingLogoColorPictureEN}

    \end{textblock}


    \begin{textblock}{2.5}(4.15,2.5)

      \includegraphics[scale=0.2]{\JULogoColorPictureEN}

    \end{textblock}


    \begin{textblock}{2.5}(6.35,2.4)

      \includegraphics[scale=0.1]{\ZintegrUJLogoColorPictureEN}

    \end{textblock}


    \begin{textblock}{4.2}(8.4,2.6)

      \includegraphics[scale=0.3]{\EUSocialFundLogoColorPictureEN}

    \end{textblock}





    \begin{textblock}{11}(1,4)

      \begin{flushleft}

        \mdseries

        \footnotesize

        \RaggedRight

        \color{jFrametitleFGColor}

        The content of this lecture is made available under a~Creative
        Commons licence (\textsc{cc}), giving the author the credits
        (\textsc{by}) and putting an obligation to share on the same terms
        (\textsc{sa}). Figures and diagrams included in the lecture are
        authored by Paweł Węgrzyn et~al., and are available under the same
        license unless indicated otherwise.\\ The presentation uses the
        Beamer Jagiellonian theme based on Matthias Vogelgesang’s
        Metropolis theme, available under license \LaTeX{} Project
        Public License~1.3c at: \colorhref{https://github.com/matze/mtheme}
        {https://github.com/matze/mtheme}.

        Typographic design: Iwona Grabska-Gradzińska \\
        \LaTeX{} Typesetting: Kamil Ziemian \\
        Proofreading: Wojciech Palacz,
        Monika Stawicka \\
        3D Models: Dariusz Frymus, Kamil Nowakowski \\
        Figures and charts: Kamil Ziemian, Paweł Węgrzyn, Wojciech Palacz

      \end{flushleft}

    \end{textblock}

  \end{frame}
}



\newcommand{\GeometryThreeDTwoSpecialEndingSlidesEN}[1]{%
  \begin{frame}[standout]


    \begin{textblock}{11}(1,0.7)

      \begin{flushleft}

        \mdseries

        \footnotesize

        \color{jFrametitleFGColor}

        This content was created as part of a project co-financed by the
        European Union within the framework of the European Social Fund
        POWR.03.05.00-00-Z309/17-00.

      \end{flushleft}

    \end{textblock}





    \begin{textblock}{10}(0,2.2)

      \tikz \fill[color=jBackgroundStyleLight] (0,0) rectangle (12.8,-1.5);

    \end{textblock}


    \begin{textblock}{3.2}(0.7,2.45)

      \includegraphics[scale=0.3]{\FundingLogoColorPictureEN}

    \end{textblock}


    \begin{textblock}{2.5}(4.15,2.5)

      \includegraphics[scale=0.2]{\JULogoColorPictureEN}

    \end{textblock}


    \begin{textblock}{2.5}(6.35,2.4)

      \includegraphics[scale=0.1]{\ZintegrUJLogoColorPictureEN}

    \end{textblock}


    \begin{textblock}{4.2}(8.4,2.6)

      \includegraphics[scale=0.3]{\EUSocialFundLogoColorPictureEN}

    \end{textblock}





    \begin{textblock}{11}(1,4)

      \begin{flushleft}

        \mdseries

        \footnotesize

        \RaggedRight

        \color{jFrametitleFGColor}

        The content of this lecture is made available under a~Creative
        Commons licence (\textsc{cc}), giving the author the credits
        (\textsc{by}) and putting an obligation to share on the same terms
        (\textsc{sa}). Figures and diagrams included in the lecture are
        authored by Paweł Węgrzyn et~al., and are available under the same
        license unless indicated otherwise.\\ The presentation uses the
        Beamer Jagiellonian theme based on Matthias Vogelgesang’s
        Metropolis theme, available under license \LaTeX{} Project
        Public License~1.3c at: \colorhref{https://github.com/matze/mtheme}
        {https://github.com/matze/mtheme}.

        Typographic design: Iwona Grabska-Gradzińska \\
        \LaTeX{} Typesetting: Kamil Ziemian \\
        Proofreading: Wojciech Palacz,
        Monika Stawicka \\
        3D Models: Dariusz Frymus, Kamil Nowakowski \\
        Figures and charts: Kamil Ziemian, Paweł Węgrzyn, Wojciech Palacz

      \end{flushleft}

    \end{textblock}

  \end{frame}





  \begin{frame}[standout]

    \begingroup

    \color{jFrametitleFGColor}

    #1

    \endgroup

  \end{frame}
}



\newcommand{\GeometryThreeDSpecialEndingSlideVideoVerOneEN}{%
  \begin{frame}[standout]

    \begin{textblock}{11}(1,0.7)

      \begin{flushleft}

        \mdseries

        \footnotesize

        \color{jFrametitleFGColor}

        This content was created as part of a project co-financed by the
        European Union within the framework of the European Social Fund
        POWR.03.05.00-00-Z309/17-00.

      \end{flushleft}

    \end{textblock}





    \begin{textblock}{10}(0,2.2)

      \tikz \fill[color=jBackgroundStyleLight] (0,0) rectangle (12.8,-1.5);

    \end{textblock}


    \begin{textblock}{3.2}(0.7,2.45)

      \includegraphics[scale=0.3]{\FundingLogoColorPictureEN}

    \end{textblock}


    \begin{textblock}{2.5}(4.15,2.5)

      \includegraphics[scale=0.2]{\JULogoColorPictureEN}

    \end{textblock}


    \begin{textblock}{2.5}(6.35,2.4)

      \includegraphics[scale=0.1]{\ZintegrUJLogoColorPictureEN}

    \end{textblock}


    \begin{textblock}{4.2}(8.4,2.6)

      \includegraphics[scale=0.3]{\EUSocialFundLogoColorPictureEN}

    \end{textblock}





    \begin{textblock}{11}(1,4)

      \begin{flushleft}

        \mdseries

        \footnotesize

        \RaggedRight

        \color{jFrametitleFGColor}

        The content of this lecture is made available under a Creative
        Commons licence (\textsc{cc}), giving the author the credits
        (\textsc{by}) and putting an obligation to share on the same terms
        (\textsc{sa}). Figures and diagrams included in the lecture are
        authored by Paweł Węgrzyn et~al., and are available under the same
        license unless indicated otherwise.\\ The presentation uses the
        Beamer Jagiellonian theme based on Matthias Vogelgesang’s
        Metropolis theme, available under license \LaTeX{} Project
        Public License~1.3c at: \colorhref{https://github.com/matze/mtheme}
        {https://github.com/matze/mtheme}.

        Typographic design: Iwona Grabska-Gradzińska;
        \LaTeX{} Typesetting: Kamil Ziemian \\
        Proofreading: Wojciech Palacz,
        Monika Stawicka \\
        3D Models: Dariusz Frymus, Kamil Nowakowski \\
        Figures and charts: Kamil Ziemian, Paweł Węgrzyn, Wojciech
        Palacz \\
        Film editing: Agencja Filmowa Film \& Television Production~--
        Zbigniew Masklak

      \end{flushleft}

    \end{textblock}

  \end{frame}
}



\newcommand{\GeometryThreeDSpecialEndingSlideVideoVerTwoEN}{%
  \begin{frame}[standout]

    \begin{textblock}{11}(1,0.7)

      \begin{flushleft}

        \mdseries

        \footnotesize

        \color{jFrametitleFGColor}

        This content was created as part of a project co-financed by the
        European Union within the framework of the European Social Fund
        POWR.03.05.00-00-Z309/17-00.

      \end{flushleft}

    \end{textblock}





    \begin{textblock}{10}(0,2.2)

      \tikz \fill[color=jBackgroundStyleLight] (0,0) rectangle (12.8,-1.5);

    \end{textblock}


    \begin{textblock}{3.2}(0.7,2.45)

      \includegraphics[scale=0.3]{\FundingLogoColorPictureEN}

    \end{textblock}


    \begin{textblock}{2.5}(4.15,2.5)

      \includegraphics[scale=0.2]{\JULogoColorPictureEN}

    \end{textblock}


    \begin{textblock}{2.5}(6.35,2.4)

      \includegraphics[scale=0.1]{\ZintegrUJLogoColorPictureEN}

    \end{textblock}


    \begin{textblock}{4.2}(8.4,2.6)

      \includegraphics[scale=0.3]{\EUSocialFundLogoColorPictureEN}

    \end{textblock}





    \begin{textblock}{11}(1,4)

      \begin{flushleft}

        \mdseries

        \footnotesize

        \RaggedRight

        \color{jFrametitleFGColor}

        The content of this lecture is made available under a Creative
        Commons licence (\textsc{cc}), giving the author the credits
        (\textsc{by}) and putting an obligation to share on the same terms
        (\textsc{sa}). Figures and diagrams included in the lecture are
        authored by Paweł Węgrzyn et~al., and are available under the same
        license unless indicated otherwise.\\ The presentation uses the
        Beamer Jagiellonian theme based on Matthias Vogelgesang’s
        Metropolis theme, available under license \LaTeX{} Project
        Public License~1.3c at: \colorhref{https://github.com/matze/mtheme}
        {https://github.com/matze/mtheme}.

        Typographic design: Iwona Grabska-Gradzińska;
        \LaTeX{} Typesetting: Kamil Ziemian \\
        Proofreading: Wojciech Palacz,
        Monika Stawicka \\
        3D Models: Dariusz Frymus, Kamil Nowakowski \\
        Figures and charts: Kamil Ziemian, Paweł Węgrzyn, Wojciech
        Palacz \\
        Film editing: IMAVI -- Joanna Kozakiewicz, Krzysztof Magda, Nikodem
        Frodyma

      \end{flushleft}

    \end{textblock}

  \end{frame}
}



\newcommand{\GeometryThreeDSpecialEndingSlideVideoVerThreeEN}{%
  \begin{frame}[standout]

    \begin{textblock}{11}(1,0.7)

      \begin{flushleft}

        \mdseries

        \footnotesize

        \color{jFrametitleFGColor}

        This content was created as part of a project co-financed by the
        European Union within the framework of the European Social Fund
        POWR.03.05.00-00-Z309/17-00.

      \end{flushleft}

    \end{textblock}





    \begin{textblock}{10}(0,2.2)

      \tikz \fill[color=jBackgroundStyleLight] (0,0) rectangle (12.8,-1.5);

    \end{textblock}


    \begin{textblock}{3.2}(0.7,2.45)

      \includegraphics[scale=0.3]{\FundingLogoColorPictureEN}

    \end{textblock}


    \begin{textblock}{2.5}(4.15,2.5)

      \includegraphics[scale=0.2]{\JULogoColorPictureEN}

    \end{textblock}


    \begin{textblock}{2.5}(6.35,2.4)

      \includegraphics[scale=0.1]{\ZintegrUJLogoColorPictureEN}

    \end{textblock}


    \begin{textblock}{4.2}(8.4,2.6)

      \includegraphics[scale=0.3]{\EUSocialFundLogoColorPictureEN}

    \end{textblock}





    \begin{textblock}{11}(1,4)

      \begin{flushleft}

        \mdseries

        \footnotesize

        \RaggedRight

        \color{jFrametitleFGColor}

        The content of this lecture is made available under a Creative
        Commons licence (\textsc{cc}), giving the author the credits
        (\textsc{by}) and putting an obligation to share on the same terms
        (\textsc{sa}). Figures and diagrams included in the lecture are
        authored by Paweł Węgrzyn et~al., and are available under the same
        license unless indicated otherwise.\\ The presentation uses the
        Beamer Jagiellonian theme based on Matthias Vogelgesang’s
        Metropolis theme, available under license \LaTeX{} Project
        Public License~1.3c at: \colorhref{https://github.com/matze/mtheme}
        {https://github.com/matze/mtheme}.

        Typographic design: Iwona Grabska-Gradzińska;
        \LaTeX{} Typesetting: Kamil Ziemian \\
        Proofreading: Wojciech Palacz,
        Monika Stawicka \\
        3D Models: Dariusz Frymus, Kamil Nowakowski \\
        Figures and charts: Kamil Ziemian, Paweł Węgrzyn, Wojciech
        Palacz \\
        Film editing: Agencja Filmowa Film \& Television Production~--
        Zbigniew Masklak \\
        Film editing: IMAVI -- Joanna Kozakiewicz, Krzysztof Magda, Nikodem
        Frodyma

      \end{flushleft}

    \end{textblock}

  \end{frame}
}



\newcommand{\GeometryThreeDTwoSpecialEndingSlidesVideoVerOneEN}[1]{%
  \begin{frame}[standout]

    \begin{textblock}{11}(1,0.7)

      \begin{flushleft}

        \mdseries

        \footnotesize

        \color{jFrametitleFGColor}

        This content was created as part of a project co-financed by the
        European Union within the framework of the European Social Fund
        POWR.03.05.00-00-Z309/17-00.

      \end{flushleft}

    \end{textblock}





    \begin{textblock}{10}(0,2.2)

      \tikz \fill[color=jBackgroundStyleLight] (0,0) rectangle (12.8,-1.5);

    \end{textblock}


    \begin{textblock}{3.2}(0.7,2.45)

      \includegraphics[scale=0.3]{\FundingLogoColorPictureEN}

    \end{textblock}


    \begin{textblock}{2.5}(4.15,2.5)

      \includegraphics[scale=0.2]{\JULogoColorPictureEN}

    \end{textblock}


    \begin{textblock}{2.5}(6.35,2.4)

      \includegraphics[scale=0.1]{\ZintegrUJLogoColorPictureEN}

    \end{textblock}


    \begin{textblock}{4.2}(8.4,2.6)

      \includegraphics[scale=0.3]{\EUSocialFundLogoColorPictureEN}

    \end{textblock}





    \begin{textblock}{11}(1,4)

      \begin{flushleft}

        \mdseries

        \footnotesize

        \RaggedRight

        \color{jFrametitleFGColor}

        The content of this lecture is made available under a Creative
        Commons licence (\textsc{cc}), giving the author the credits
        (\textsc{by}) and putting an obligation to share on the same terms
        (\textsc{sa}). Figures and diagrams included in the lecture are
        authored by Paweł Węgrzyn et~al., and are available under the same
        license unless indicated otherwise.\\ The presentation uses the
        Beamer Jagiellonian theme based on Matthias Vogelgesang’s
        Metropolis theme, available under license \LaTeX{} Project
        Public License~1.3c at: \colorhref{https://github.com/matze/mtheme}
        {https://github.com/matze/mtheme}.

        Typographic design: Iwona Grabska-Gradzińska;
        \LaTeX{} Typesetting: Kamil Ziemian \\
        Proofreading: Wojciech Palacz,
        Monika Stawicka \\
        3D Models: Dariusz Frymus, Kamil Nowakowski \\
        Figures and charts: Kamil Ziemian, Paweł Węgrzyn,
        Wojciech Palacz \\
        Film editing: Agencja Filmowa Film \& Television Production~--
        Zbigniew Masklak

      \end{flushleft}

    \end{textblock}

  \end{frame}





  \begin{frame}[standout]


    \begingroup

    \color{jFrametitleFGColor}

    #1

    \endgroup

  \end{frame}
}



\newcommand{\GeometryThreeDTwoSpecialEndingSlidesVideoVerTwoEN}[1]{%
  \begin{frame}[standout]

    \begin{textblock}{11}(1,0.7)

      \begin{flushleft}

        \mdseries

        \footnotesize

        \color{jFrametitleFGColor}

        This content was created as part of a project co-financed by the
        European Union within the framework of the European Social Fund
        POWR.03.05.00-00-Z309/17-00.

      \end{flushleft}

    \end{textblock}





    \begin{textblock}{10}(0,2.2)

      \tikz \fill[color=jBackgroundStyleLight] (0,0) rectangle (12.8,-1.5);

    \end{textblock}


    \begin{textblock}{3.2}(0.7,2.45)

      \includegraphics[scale=0.3]{\FundingLogoColorPictureEN}

    \end{textblock}


    \begin{textblock}{2.5}(4.15,2.5)

      \includegraphics[scale=0.2]{\JULogoColorPictureEN}

    \end{textblock}


    \begin{textblock}{2.5}(6.35,2.4)

      \includegraphics[scale=0.1]{\ZintegrUJLogoColorPictureEN}

    \end{textblock}


    \begin{textblock}{4.2}(8.4,2.6)

      \includegraphics[scale=0.3]{\EUSocialFundLogoColorPictureEN}

    \end{textblock}





    \begin{textblock}{11}(1,4)

      \begin{flushleft}

        \mdseries

        \footnotesize

        \RaggedRight

        \color{jFrametitleFGColor}

        The content of this lecture is made available under a Creative
        Commons licence (\textsc{cc}), giving the author the credits
        (\textsc{by}) and putting an obligation to share on the same terms
        (\textsc{sa}). Figures and diagrams included in the lecture are
        authored by Paweł Węgrzyn et~al., and are available under the same
        license unless indicated otherwise.\\ The presentation uses the
        Beamer Jagiellonian theme based on Matthias Vogelgesang’s
        Metropolis theme, available under license \LaTeX{} Project
        Public License~1.3c at: \colorhref{https://github.com/matze/mtheme}
        {https://github.com/matze/mtheme}.

        Typographic design: Iwona Grabska-Gradzińska;
        \LaTeX{} Typesetting: Kamil Ziemian \\
        Proofreading: Wojciech Palacz,
        Monika Stawicka \\
        3D Models: Dariusz Frymus, Kamil Nowakowski \\
        Figures and charts: Kamil Ziemian, Paweł Węgrzyn,
        Wojciech Palacz \\
        Film editing: IMAVI -- Joanna Kozakiewicz, Krzysztof Magda, Nikodem
        Frodyma

      \end{flushleft}

    \end{textblock}

  \end{frame}





  \begin{frame}[standout]


    \begingroup

    \color{jFrametitleFGColor}

    #1

    \endgroup

  \end{frame}
}



\newcommand{\GeometryThreeDTwoSpecialEndingSlidesVideoVerThreeEN}[1]{%
  \begin{frame}[standout]

    \begin{textblock}{11}(1,0.7)

      \begin{flushleft}

        \mdseries

        \footnotesize

        \color{jFrametitleFGColor}

        This content was created as part of a project co-financed by the
        European Union within the framework of the European Social Fund
        POWR.03.05.00-00-Z309/17-00.

      \end{flushleft}

    \end{textblock}





    \begin{textblock}{10}(0,2.2)

      \tikz \fill[color=jBackgroundStyleLight] (0,0) rectangle (12.8,-1.5);

    \end{textblock}


    \begin{textblock}{3.2}(0.7,2.45)

      \includegraphics[scale=0.3]{\FundingLogoColorPictureEN}

    \end{textblock}


    \begin{textblock}{2.5}(4.15,2.5)

      \includegraphics[scale=0.2]{\JULogoColorPictureEN}

    \end{textblock}


    \begin{textblock}{2.5}(6.35,2.4)

      \includegraphics[scale=0.1]{\ZintegrUJLogoColorPictureEN}

    \end{textblock}


    \begin{textblock}{4.2}(8.4,2.6)

      \includegraphics[scale=0.3]{\EUSocialFundLogoColorPictureEN}

    \end{textblock}





    \begin{textblock}{11}(1,4)

      \begin{flushleft}

        \mdseries

        \footnotesize

        \RaggedRight

        \color{jFrametitleFGColor}

        The content of this lecture is made available under a Creative
        Commons licence (\textsc{cc}), giving the author the credits
        (\textsc{by}) and putting an obligation to share on the same terms
        (\textsc{sa}). Figures and diagrams included in the lecture are
        authored by Paweł Węgrzyn et~al., and are available under the same
        license unless indicated otherwise. \\ The presentation uses the
        Beamer Jagiellonian theme based on Matthias Vogelgesang’s
        Metropolis theme, available under license \LaTeX{} Project
        Public License~1.3c at: \colorhref{https://github.com/matze/mtheme}
        {https://github.com/matze/mtheme}.

        Typographic design: Iwona Grabska-Gradzińska;
        \LaTeX{} Typesetting: Kamil Ziemian \\
        Proofreading: Leszek Hadasz, Wojciech Palacz,
        Monika Stawicka \\
        3D Models: Dariusz Frymus, Kamil Nowakowski \\
        Figures and charts: Kamil Ziemian, Paweł Węgrzyn,
        Wojciech Palacz \\
        Film editing: Agencja Filmowa Film \& Television Production~--
        Zbigniew Masklak \\
        Film editing: IMAVI -- Joanna Kozakiewicz, Krzysztof Magda, Nikodem
        Frodyma


      \end{flushleft}

    \end{textblock}

  \end{frame}





  \begin{frame}[standout]


    \begingroup

    \color{jFrametitleFGColor}

    #1

    \endgroup

  \end{frame}
}











% ------------------------------------------------------
% BibLaTeX
% ------------------------------------------------------
% Package biblatex, with biber as its backend, allow us to handle
% bibliography entries that use Unicode symbols outside ASCII.
\usepackage[
language=polish,
backend=biber,
style=alphabetic,
url=false,
eprint=true,
]{biblatex}

\addbibresource{Systemy-operacyjne-Bibliography.bib}





% ------------------------------------------------------
% Importing packages, libraries and setting their configuration
% ------------------------------------------------------





% ------------------------------------------------------
% Local packages
% ------------------------------------------------------
% Local configuration of this particular presentation
\usepackage{./Local-packages/local-settings}










% ####################################################################
\title{Systemy operacyjne}
\subtitle{XX. Procesy w~systemach operacyjnych}

\author{Kamil Ziemian \\
  \email}


% \date{}
% ####################################################################










% ####################################################################
% Beginning of the document
\begin{document}
% ####################################################################





% ################################################
% Number of chars: 20k+, 26k+, 41k+,
% Text is adjusted to the left and words are broken at the end of the line.
\RaggedRight
% ################################################





% ################################################
\maketitle
% ################################################





% ################################################
\EndingSlide{Ta prezentacja wciąż nie jest ukończona.}
% ################################################





% ################################################
\begin{frame}
  \frametitle{Spis treści}


  \tableofcontents

\end{frame}
% ################################################





% ################################################
\section{Do czego służą procesy?}
% ################################################


% ######################################
\begin{frame}
  \frametitle{Do czego służą procesy?}


  Naszą opowieść o~procesach zaczniemy od dwóch cytatów.

  \textit{Wszystko w~systemie Linux jest plikiem. Jeśli coś nie jest
    plikiem to jest procesem.} \\
  Na ten moment, %\date,
  autora tego powiedzenie nie udało~się ustalić.

  \textit{Ogólnie proces jest kontenerem zawierającym wszystkie informacje
    niezbędne do uruchomienia programu.} \\
  Andrewa S.~Tanenbaum, str.~$72$,
  \parencite{Tanenbaum-Systemy-operacyjne-Wydanie-III-Pub-2013}.

  Dodajmy, że~przez „kontener” (ang.~\textit{container}) Tanenbaum rozumie
  rodzaj pojemnika, w~który zawarte są wszystkie potrzebne nam rzeczy.

  Jak przypomina pierwszy cytat, proces jest obok pliku, jednym z~dwóch
  najważniejszych pojęć w~systemie operacyjnym GNU/Linux. Drugi cytat
  może być trudniejszy do zrozumienie, proszę~się jednak nie martwić,
  jeśli nie wiedzą Państwo co Tanenbaum chciał przekazać, choć jego słowa
  kryją głęboką treść.

\end{frame}
% ######################################





% ######################################
\begin{frame}
  \frametitle{Do czego służą procesy?}


  Ponieważ procesy są tak fundamentalne dla działania komputerów, można
  by spędzić kilka dni dyskutując czym są, czy jak działają. Natomiast
  problemowi jak zarządzać procesami trzeba by poświęcić kilka miesięcy
  czasu. Jednak nasz przedmiot jest tylko wstępem do systemów operacyjnych,
  więc zadowolimy~się bardzo \alert{uproszczonym} spojrzeniem na całą
  tematykę procesów.

  Przypomnijmy jeden ze~sloganów tego przedmiotu, który brzmi następująco.
  Plik czymś jest, proces coś robi. W~innej wersji: plik zawiera dane,
  proces przetwarza dane.

  Należy zrozumieć, że~pliki to obiekty dość bierne i~mocno statyczne.
  Z~punktu widzenia komputera większość plików zmienia~się stosunkowo
  rzadko. W~przeciwieństwo do nich procesy są aktywne, dynamiczne,
  zmieniające swój stan cały czas. Spróbujemy wyjaśnić to na kilku
  przykładach.

\end{frame}
% ######################################





% ######################################
\begin{frame}
  \frametitle{Dlaczego potrzebujemy procesów?}


  Powiedzmy, że~w~katalogu ma plik \texttt{Data-01.dat}, który chcemy
  przenieść do katalogu \texttt{Data-March/}. By to zrobić użyjemy
  polecenia \texttt{mv -i}: \\
  \texttt{\$ mv -i Data-01.dat Data-March/} \\
  Ponieważ polecenie \texttt{mv -i} wykonuje pewną czynność, więc należy
  podejrzewać, że~jego uruchomienie utworzy proces, który to proces
  przeniesie plik \texttt{Data-01.dat} do katalogu. Tak w~istocie jest.
  Choć zwykle tego nie dostrzegamy, to wywołanie \alert{prawie każdego}
  polecenie skutkuje powstanie procesu, który wykonuje daną czynność.

  Jest to tak ważne, że~warto to powiedzieć jeszcze raz, w~trochę innym
  sposób. Z~zasady, jeśli chcemy by komputery wykonał \alert{dowolną}
  czynność, to potrzebujemy pewnego procesu, który tą czynność wykona.

  To wciąż nam niewiele mówi na temat tego, jak proces jest zbudowany
  wewnętrznie, jednak na tym przedmiocie dużo ważniejsze jest zrozumienie
  na poziomie ogólnym do czego służą procesy, niż to jak są wewnętrznie
  zbudowane.

\end{frame}
% ######################################





% ######################################
\begin{frame}
  \frametitle{Dlaczego potrzebujemy procesów?}


  Rozpatrzmy teraz przykład z~innej bajki. Załóżmy, że~chcemy sprawdzić,
  kiedy przebieg pierwszych miesięcy I~Wojny Światowej, które są opisane
  w~książce historycznej, którą mamy na półce. Sięgamy tę książkę z~półki
  i~zaczynamy czytać. W~tej opowieści książka to odpowiednik pliku,
  ściągnięcie jej z~półki to odpowiednik procesu.

  Czemu ma służyć ta opowieść? Zauważmy, że~wedle niej książka zawiera
  pewne informacje/dane, ale zawiera je w~sposób bierny. Informacje w~niej
  są i~koniec. Jednak, aby móc skorzystać z~tych informacji musimy wykonać
  pewne czynności, takie jak zdjęcie książki z~półki. Gdy chodzi o~książki,
  to nie nastręcza problemu, jednak w~przypadku komputerów należy
  zwrócić uwagę na kilka rzeczy.

  Pliki ze swej natury to dość bierne obiekty. przechowują one pewne dane
  i~nic więcej. Mówiąc ogólnie, gdybyśmy mieli dany plik z~danymi, ale nie
  mogli uruchomić żadnego procesu, to dane te istniałyby, ale byłby
  zupełnie bezużyteczne, bo nie umielibyśmy ich odczytać.

\end{frame}
% ######################################





% ######################################
\begin{frame}
  \frametitle{Dlaczego potrzebujemy procesów?}


  Używając poprzedniego przykładu, pożądane informacje są  w~danej książce,
  ale jeśli nie ściągniemy jej z~półki, nigdy do nich nie dotrzemy. Gdy
  chodzi o~dane zawarte w~pliku by~je uzyskać, potrzebujemy procesu, który
  otworzy plik, pobierze z~niego dane i~wyświetli je nam na ekranie.

  Zarówno gdy chodzi o~książkę, jak i~o~dane zawarte w~pliku, to na końcu
  musimy coś przeczytać, zrozumieć znaczenie danego tekstu i~dopiero wtedy
  można powiedzieć, że~przyswoiliśmy sobie pewne informacje. W~jaki sposób
  człowiek czyta i~przyswaja informacje to jednak zagadnienie spoza
  informatyki, którego dyskusji z~premedytacją będziemy~się wystrzegać.
  Nie mamy bowiem nic mądrego do~powiedzenia w~tej kwestii.

  Historia ta ma ilustrować następującą rzecz. W~komputerze nawet
  wyświetlenie danych z~pliku wymaga procesu, bo musimy odczytać je z~dysku
  lub innego miejsca pamięci i~wyświetlić je na ekranie.

\end{frame}
% ######################################





% ######################################
\begin{frame}
  \frametitle{Do czego służą procesy?}


  Widzimy, że~nawet coś takiego jak wyświetlanie zawartości pliku jest
  formą przetwarzania danych. Zagadnienie w~jaki sposób dane są pobierane
  z~dysku i~wyświetlane na ekranie to sam w~sobie duży temat, w~który nie
  będziemy~się więcej zagłębiać, bo nie należy to stricte do tematu
  naszego przedmiotu.

  Najważniejsze dla nas jest to, że~do każdej czynności na komputerze,
  nawet tak niewinnej jak przeglądanie zawartości pliku, potrzebujemy
  jakiegoś procesu, który ją wykona.

  Dla uczciwości należy przypomnieć, że~przedstawiony opis czytania z~pliku
  jest pewnym uproszczeniem rzeczywistości, niemniej ogólny zarys problemu
  jest poprawny. Konkretną realizacją tego o~czym mówiliśmy, gdy chodzi
  o~wyświetlanie zawartości pliku, jest użycie polecenia takiego jak
  \texttt{cat}, by~wyświetlić zawartość pewnego pliku.

\end{frame}
% ######################################





% ######################################
\begin{frame}
  \frametitle{Przykłady procesów}


 W~przykładzie użyliśmy pliku
  \texttt{PiR-Część-I.txt}, dostępnego pośród innych materiałów do kursu. \\
  \texttt{\$ cat PiR-Część-I.txt} \\
  \texttt{::>} \texttt{Poniższy tekst jest zapisem podróży i~snutych
    podczas niej refleksji, jaką} \\
  \texttt{::>} \texttt{pewna dość specyficzna osoba odbyła po Polsce} \\
  \texttt{i~okolicach. Nazwa tego} \\
  \texttt{::>} \texttt{pliku „PiR-Część-I.txt” to skrót od „Podróże} \\
  \texttt{i~rozmyślania. Część I”. Oprócz} \\
  \hspace{2em} \texttt{\vdots}

  Rozpatrzmy jeszcze jeden przykład procesu. Na dysku mamy program plik
  \texttt{sum}, zawierający prosty program. Program ten pobiera
  jako argument nawę pliku i~oblicza sumę liczb jaką ten program zawierał.

\end{frame}
% ######################################





% ######################################
\begin{frame}
  \frametitle{Przykłady procesów}


  Możemy myśleć, że~program \texttt{sum} został napisany w~języku takim
  jak C, Python, czy jakimś innym. Choć w~przypadku języka Python powinien
  on nosić nazwę taką jak \texttt{sum.py}, nie musimy~się tym jednak teraz
  przejmować. Kod tego programu zapewne będzie zawierał pętlę \texttt{for}
  służącą do obliczenia sumy liczb jakie program otrzyma na~wejściu.
  Wszystkie instrukcje jakie musi on wykonać na pliku i~zawartych
  w~nim liczbach są zapisane w~pliku \texttt{sum}.

\end{frame}
% ######################################





% ######################################
\begin{frame}
  \frametitle{Korzystanie}




\end{frame}
% ######################################


% ######################################
\begin{frame}
  \frametitle{Duże i~małe litery}


  Domyślnie \texttt{grep} rozróżnia duże i~małe liter, co~obrazują poniższe
  przykłady. \\
  \texttt{\$ grep -n "Po"{} PiR-Część-I.txt} \\
  \texttt{::> 1:\emphasisRed{Po}niższy tekst, jest zapisem podroży
    i~snutych} \\
  \texttt{podczas niej refleksji, jaką} \\
  \texttt{::> 2:pewna dość specyficzna osoba odbyła
    po~\emphasisRed{Po}lsce} \\
  \texttt{i~okolicach. Nazwa tego} \\
  \hspace{2em} \texttt{\vdots}

  \texttt{\$ grep -n "po"{} PiR-Część-I.txt} \\
  \texttt{::> 1:Poniższy tekst, jest zapisem \emphasisRed{po}droży
    i~snutych} \\
  \texttt{\emphasisRed{po}dczas niej refleksji, jaką} \\
  \texttt{::> 2:pewna dość specyficzna osoba odbyła \emphasisRed{po}
    Polsce} \\
  \texttt{i~okolicach. Nazwa tego} \\
  \hspace{2em} \texttt{\vdots}

\end{frame}
% ######################################





% ######################################
\begin{frame}
  \frametitle{Duże i~małe litery}


  Jeśli chcemy, by \texttt{grep} ignorował rozmiar liter, to używamy
  opcji \texttt{-i} od ang.~\textit{ignore case}, pl.~\textit{ignoruj
    wielkość liter}. \\
  \texttt{\$ grep -n -i "Po"{} PiR-Część-I.txt} \\
  \texttt{::> 1:\emphasisRed{Po}niższy tekst, jest zapisem
    \emphasisRed{po}droży i~snutych} \\
  \texttt{\emphasisRed{po}dczas niej refleksji, jaką} \\
  \texttt{2:pewna dość specyficzna osoba odbyła \emphasisRed{po}
    \emphasisRed{Po}lsce} \\
  \texttt{i~okolicach. Nazwa tego} \\
  \hspace{2em} \texttt{\vdots}

  Równoważnie, można użyć następującego polecenia. \\
  \texttt{\$ grep -n --ignore-case "Po"{} PiR-Część-I.txt} \\
  \texttt{::> 1:\emphasisRed{Po}niższy tekst, jest zapisem
    \emphasisRed{po}droży i~snutych} \\
  \texttt{\emphasisRed{po}dczas niej refleksji, jaką} \\
  \texttt{2:pewna dość specyficzna osoba odbyła \emphasisRed{po}
    \emphasisRed{Po}lsce} \\
  \texttt{i~okolicach. Nazwa tego} \\
  \hspace{2em} \texttt{\vdots}

\end{frame}
% ######################################





% ######################################
\begin{frame}
  \frametitle{Korzystanie z~polecenia \texttt{grep}}


  Jeśli chcemy wyświetlić tylko te linie, które \alert{nie} zawierają
  danego tekstu, można skorzystać z~opcji~\texttt{-v}. By to zilustrować,
  porównamy rezultaty użycia tych dwóch wersji tego polecenia. \\
  \texttt{\$ grep -n -v "Po"{} PiR-Część-I.txt} \\
  \texttt{::> 4:tego, że bohater ma skłonność do opisywania tego co widzi
    w~terminach nauki,} \\
  \texttt{::> 5:to snuje refleksje o muzyce, literaturze i~wielu innych
    rzeczach, które} \\
  \texttt{::> 6:przyjdą mu do głowy.} \\
  \texttt{::> 7:} \\
  \texttt{::> 8:Do Lublina dojechałem samochodem o~godzinie 18:30} \\
  \texttt{i~skierowałem~się} \\
  \hspace{2em} \texttt{\vdots}

\end{frame}
% ######################################





% ######################################
\begin{frame}
  \frametitle{Jakie tekstu może szukać \texttt{grep}?}


  Warto zwrócić uwagę, że~polecenie \texttt{grep} na poprzednim slajdzie
  zwróciło też linię numer~$7$, która jest zupełnie pusta. Jest to poprawne
  zachowanie z~jej strony, jako~że pusta linie rzeczywiście nie zawiera
  ciągu znaków~„Po”.

  Do tej pory używaliśmy \texttt{grep}a tylko do wyszukiwania ciągów
  liter, jednak jego możliwości nie ograniczają~się tylko do takich
  zastosowań. Równie dobrze możemy wyszukiwać ciągi znaków
  zawierające spację, co ilustruje poniższy przykład. \\
  \texttt{\$ grep -n "a p"{} PiR-Część-I.txt} \\
  \texttt{::> 2:pewna dość specyficzna osob\emphasisRed{a o}dbyła
    po~Polsce} \\
  \texttt{i~okolicach. Nazwa tego} \\
  \texttt{::> 31:obejrzałem. Z~tego też powodu zapewne masa ludzi} \\
  \texttt{n\emphasisRed{a p}ewno go nie znosi, bo} \\
  \hspace{2em} \texttt{\vdots}

\end{frame}
% ######################################





% ######################################
\begin{frame}
  \frametitle{Jakie tekstu może szukać \texttt{grep}?}


  Tak samo możemy poszukiwać w~tekście cyfry. By wyszukać wystąpienie
  w~tekście cyfry „$1$” można użyć poniższej wersji polecenia
  \texttt{grep}. \\
  \texttt{\$ grep -n "1"{} PiR-Część-I.txt} \\
  \texttt{::> 8:Do Lublina dojechałem samochodem o~godzinie
    \emphasisRed{1}8:30}
  \texttt{i~skierowałem~się} \\
  \texttt{12:mam pierwszy tom jego opowiadań zebranych, wydany} \\
  \texttt{w~Polsce w~\emphasisRed{1}967 roku.} \\
  \hspace{2em} \texttt{\vdots}

  Oczywiście możemy też szukać tekstu, który składa~się z~liter, spacji,
  cyfr,~etc., jak w~poniższym przykładzie. \\
  \texttt{\$ grep -n "e 1"{} PiR-Część-I.txt} \\
  \texttt{::> 8:Do Lublina dojechałem samochodem}
  \texttt{o~godzini\emphasisRed{e 1}8:30 i~skierowałem~się}

\end{frame}
% ######################################





% ######################################
\begin{frame}
  \frametitle{Jakie tekstu może szukać \texttt{grep}?}


  W~tym wypadku tylko jedna linia zawiera szukany przez nas tekst. Powinno
  teraz być jasne, że~\texttt{grep}owi można podać w~zasadzie dowolny ciąg
  symboli i~on znajdzie wszystkie linie w~których~się on~pojawia.

  Poleceniu \texttt{grep} można przesłać w~zasadzie dowolny tekst,
  jaki tylko jesteśmy w~stanie wymyślić, jest tu tylko jedno poważne
  zastrzeżenie. Mianowicie, \texttt{grep} zakłada, że~podany mu tekst
  zawiera~się cały w~\alert{pojedynczej linii}. Zobaczymy jak to działa na
  przykładzie.

  W~pliku \texttt{PiR-Część-I.txt} znajdują~się następujące dwie linie. \\
  \textit{3:pliku „PiR-Część-I.txt” to skrót od „Podróże i rozmyślania.
    Część~I”. Oprócz} \\
  \textit{4:tego, że bohater ma skłonność do opisywania tego
    co~widzi w~terminach nauki,} \\
  Dla wygody do każdej z~nich dodaliśmy jej numer. Jak widzimy, tekst ten
  zawiera ciąg słów „Oprócz tego”, ale rozbity na dwie linie, wobec
  czego \texttt{grep} go nie znajdzie.

\end{frame}
% ######################################





% ######################################
\begin{frame}
  \frametitle{Jakie tekstu może szukać \texttt{grep}?}


  By~to zobaczyć, proszę~się przyjrzeć następującym przykładom. \\
  \texttt{\$ grep -n "Oprócz"{} PiR-Część-I.txt} \\
  \texttt{::> 3:pliku „PiR-Część-I.txt” to skrót od „Podróże} \\
  \texttt{i~rozmyślania. Część I”. \emphasisRed{Oprócz}}

  \texttt{\$ grep -n "tego"{} PiR-Część-I.txt} \\
  \texttt{::> 2:pewna dość specyficzna osoba odbyła po~Polsce} \\
  \texttt{i~okolicach. Nazwa \emphasisRed{tego}} \\
  \texttt{::> 4:\emphasisRed{tego}, że~bohater ma skłonność do opisywania
    \emphasisRed{tego}} \\
  \texttt{co widzi w~terminach nauki,} \\
  \hspace{2em} \texttt{\vdots}

  Jak należało~się spodziewać, polecenie \texttt{grep} nie ma problemu
  ze~znalezieniem słowa „Oprócz” i~„tego”,

\end{frame}
% ######################################





% ######################################
\begin{frame}
  \frametitle{Zliczanie znalezionych linii}


  Teraz poszukajmy tekstu „Oprócz tego”. \\
  \texttt{\$ grep -n "Oprócz tego"{} PiR-Część-I.txt} \\
  \texttt{::>} \\
  Jak Państwo widzą, \texttt{grep} nie znalazł tekstu „Oprócz tego”, gdyż
  ten jest podzielony na dwie linie.

  By~polecenie \texttt{grep} wypisało w~jakiej liczbie linii został
  znaleziony zadany ciąg symboli (do ilu linii został dopasowany wzorzec)
  dajemy mu flagę~\texttt{-c}: \\
  \texttt{\$ grep -n -c "nie"{} PiR-Część-I.txt} \\
  \texttt{::>} \texttt{47}

  Zobaczmy w~ilu liniach „nie” występuje jako osobne słowo. \\
  \texttt{\$ grep -n -w -c "nie"{} PiR-Część-I.txt} \\
  \texttt{::>} \texttt{18}

\end{frame}
% ######################################





% ######################################
\begin{frame}
  \frametitle{Zliczanie znalezionych linii}


  Dla porządku warto zaznaczyć, że~gdy szukamy w~ilu liniach „nie”
  występuje jako osobne słowo, opcja \texttt{-n} nie jest potrzebna.
  Co więcej wszystkie poniższe wersje polecenia \texttt{grep} są
  równoważne. \\
  \texttt{grep -w -c "nie"{} PiR-Część-I.txt} \\
  \texttt{grep -c -w "nie"{} PiR-Część-I.txt} \\
  \texttt{grep -wc "nie"{} PiR-Część-I.txt} \\
  \texttt{grep -cw "nie"{} PiR-Część-I.txt}

  Należy podkreślić, że~\texttt{grep} z~opcją~\texttt{-c} zlicza ile
  \alert{linii} zawiera dany tekst, nie ile razy ten tekst wystąpił.
  Jeśli by więc tekst „nie” występował tylko w~jednej linii, ale~trzy
  razy, to~polecenie \texttt{grep -c} zwróci jako rezultat liczbę~$1$.

\end{frame}
% ######################################





% ######################################
\begin{frame}
  \frametitle{Wyszukiwanie linii w~których nie ma danego tekstu}


  Jak już wiemy, domyślnie \texttt{grep} wypisuje linie w~których znalazł
  dany ciąg symboli. By~zobaczyć linie \alert{nie} zawierające tego ciągu
  należy użyć tego polecania z~opcją \texttt{-v}. \\
  \texttt{\$ grep -n -v "nie"{} PiR-Część-I.txt} \\
  \texttt{::>} \texttt{2:pewna dość specyficzna osoba odbyła po Polsce} \\
  \texttt{i~okolicach. Nazwa tego} \\
  \texttt{::>} \texttt{3:pliku „PiR-Część-I.txt” to skrót od „Podróże} \\
  \texttt{i~rozmyślania. Część I”. Oprócz} \\
  \hspace{2em} \texttt{\vdots}

  Powiedzmy, że~oprócz linii w~której znaleziono ciąg symboli, chcemy
  zobaczyć też linię która następuje zaraz za~nią. W~tym celu do polecenia
  dodajemy opcję \texttt{-A 1}. Analogicznie, by zobaczyć dwie linie
  następujące po każdej linii w~której znaleziono ten ciąg dodajemy
  flagę \texttt{-A 2}, gdy chcemy zobaczyć trzy linie flagę
  \texttt{-A 3}, etc.

\end{frame}
% ######################################





% ######################################
\begin{frame}
  \frametitle{Jak zobaczyć szerszy kontekst?}


  \texttt{\$ grep -n -A 1 "Faulkner"{} PiR-Część-I.txt} \\
  \texttt{::>} \texttt{11:gdzie chciałem dalej czytać tom opowiadań
    Williama \emphasisRed{Faulkner}a. Dokładniej,} \\
  \texttt{::>} \texttt{12-mam pierwszy tom jego opowiadań zebranych,
    wydany} \\
  \texttt{w~Polsce w~1967 roku.} \\
  \texttt{::>} \texttt{--} \\
  \texttt{::>} \texttt{16:dobrze napisane. \emphasisRed{Faulkner} to
    niewątpliwie bardzo dobry pisarz.} \\
  \texttt{::>} \texttt{17-} \\
  \hspace{2em} \texttt{\vdots}

  Symbol „\texttt{--}” oznacza, że~kończy~się kontekst związany z~jednym
  wystąpienia szukanego ciągu znaków, w~naszym wypadku jest nim „Faulkner”.
  Następna linia po symbolu „\texttt{--}” zawiera kontekst dotyczący
  kolejnego wystąpienia szukanego ciągu znaków.

\end{frame}
% ######################################





% ######################################
\begin{frame}
  \frametitle{Jak zobaczyć szerszy kontekst?}


  \texttt{\$ grep -n -A 2 "Faulkner"{} PiR-Część-I.txt} \\
  \texttt{::>} \texttt{11:gdzie chciałem dalej czytać tom opowiadań
    Williama \emphasisRed{Faulkner}a. Dokładniej,} \\
  \texttt{::>} \texttt{12-mam pierwszy tom jego opowiadań zebranych, wydany
    w~Polsce w~1967 roku.} \\
  \texttt{::>} \texttt{13-Otwiera go opowiadanie „Spalenie stodoły”, które
    jest tak mocno osadzone} \\
  \texttt{--} \\
  \texttt{::>} \texttt{16:dobrze napisane. Faulkner to niewątpliwie bardzo
    dobry pisarz.} \\
  \hspace{2em} \texttt{\vdots}

  Symbol „\texttt{-}” w~„\texttt{12-}” oznacza, że~w~linii numer~$12$ nie
  znaleziono danego ciągu symboli i~jest ona wyświetlana dla kontekstu.

\end{frame}
% ######################################





% ######################################
\begin{frame}
  \frametitle{Jak zobaczyć szerszy kontekst?}


  Aby zobaczyć wiersz, który znajduje~się przed tym w~którym znaleziono
  szukany ciąg symboli, używamy opcję \texttt{-B 1}. Analogicznie działają
  opcję \texttt{-B 2}, \texttt{-B 3}, etc.

  \texttt{\$ grep -n -B 1 "Faulkner"{} PiR-Część-I.txt} \\
  \texttt{::>} \texttt{10-rzadka rzecz w hotelach. Po zjedzeniu kolacji} \\
  \texttt{udałem~się do swojego pokoju,} \\
  \texttt{::>} \texttt{11:gdzie chciałem dalej czytać tom opowiadań
    Williama \emphasisRed{Faulkner}a. Dokładniej,} \\
  \texttt{::>} \texttt{--} \\
  \texttt{::>} \texttt{15-i trudne do zrozumienia, pomimo tego, iż samo} \\
  \texttt{opowiadanie jest naprawdę} \\
  \hspace{2em} \texttt{\vdots}

  Znaczenie symbolu „\texttt{10-}” jest takie same jak poprzednio.

\end{frame}
% ######################################





% ######################################
\begin{frame}
  \frametitle{Jak zobaczyć szerszy kontekst?}




\end{frame}
% ######################################



% ######################################
\begin{frame}
  \frametitle{Do czego służą procesy?}




\end{frame}
% ######################################





% ##################
\begin{frame}
  \frametitle{Czym jest proces?}


  Ze~względu na to, że~produkcja takiej szybkiej pamięci jest bardzo droga,
  rejestry są zwykle bardzo małe, zawierają często tylko \alert{kilka}
  bajtów. Problem produkcji szybkiej pamięci, która jest jednocześnie
  odpowiednio tania, to ciągła bolączka elektroniki XXI~wieku. Przejdźmy
  teraz do procesów.

   Powiedzenie popularne w~środowisku
  użytkowników tego systemu, wyjaśniające jedną z~głównych zasad
  jego funkcjonowania.

  Jeśli
  komputer z~systemem GNU/Linux coś robi, to odpowiada za~to jeden lub
  więcej procesów. Innym określeniem procesu, jest stwierdzenie, że~jest
  to \alert{działający program}.

\end{frame}
% ##################





% ##################
\begin{frame}
  \frametitle{Zaglądnijmy pod warstwę abstrakcji procesu}


  Czy jest jakaś różnica między pojęciem procesu, a~pojęciem programu?
  To jak używamy tych terminów wskazuje na to, że~tak różnica istotnie
  istnieje. Rozpatrzmy przykładowo grę \textit{Hollow Knight} (program
  równie dobry jak każdy inny), którą mam zainstalowaną na dysku, ale
  w~którą akurat teraz nie gram. Gra ta niewątpliwie jest programem, ale
  nie jest nią związany żaden proces. Program więc może być więc istnieć
  w~stanie biernym, jako zbiór plików na dysku, ale nie posiadać w~danej
  chwili żadnego procesu.

  Wewnętrzna budowa procesu nie jest wcale rzeczą banalną. Każdy proces,
  z~definicji, coś robi, a~na komputerze robienie czegoś jest w~zasadzie
  równoważne przetwarzaniu (przeliczaniu) pewnego zbioru danych. Dane te
  muszą znajdować~się gdzieś w~pamięci komputera, więc każdy proces musi
  posiadać przyznany mu specjalny obszar pamięci. Ten fragment pamięci
  należący do danego procesu nazywamy \textbf{przestrzenią adresową
    procesu}.

\end{frame}
% ##################





% ##################
\begin{frame}
  \frametitle{Zaglądnijmy pod warstwę abstrakcji procesu}


  Przestrzeń adresowa danego procesu zawiera skopiowany z~dysku kod programu
  który będzie dany proces wykonywał, dane na których będzie on operował
  oraz
  \colorhref{https://en.wikipedia.org/wiki/Stack_(abstract_data_type)}{stos}
  tego procesu. W~systemach z~rodziny \textsc{unix}a część pamięci
  zawierającej sam program nazywa~się \textbf{segmentem tekstu},
  ta zawierające dane na których będzie operował \textbf{segmentem danych},
  zaś ostatnia \textbf{segmentem stosu}. Ze względu na charakter kursu nie
  możemy~się bardziej zagłębiać w~tą tematykę.

  Oprócz tego proces korzysta często z~zewnętrzny wobec niego zasobów,
  takich jak wartości które mają być wpisane do rejestrów urządzeń,
  z~których korzysta dany proces, listę otwartych plików, etc. Przykładowo
  program do wyświetlania zdjęć, by wyświetlić konkretne zdjęcie musi
  otworzyć plik je zawierające oraz ustawić rejestry monitora w~taki
  sposób, by monitor wyświetlił to zdjęcie w~odpowiednim miejscu ekranu.

\end{frame}
% ##################





% ##################
\begin{frame}
  \frametitle{Zaglądnijmy pod warstwę abstrakcji procesu}


  Dane dotyczące zewnętrznych zasobów \alert{nie} są przechowywane
  w~przestrzeni adresowej danego procesu, lecz zarządza nimi sam system
  operacyjny. W~systemie GNU/Linux jednym z~głównych narzędzi
  służących do tego jest tablica procesów.

  \textbf{Tablica procesów} jest tablicą używaną przez system operacyjny,
  w~której w razie potrzeby zapisuje on wszystkie informacje potrzebne
  do~działania danego procesu, a~które to informacje nie są zawarte w~jego
  przestrzeni adresowej.



\end{frame}
% ##################





% ##################
\begin{frame}
  \frametitle{Jaki stąd morał?}


  To co powinni Państwo wynieść z~tych uwag, to fakt, że~proces jest dość
  złożoną abstrakcją. W~dalszym ciągu zajęć nie będziemy~się zagłębiać
  w~ten temat, skupimy~się za to na tym, jak w~praktyce pracować
  z~procesami. Część teoretyczną ograniczy w~miarę możliwości do tych
  pojęć, które są niezbędne w~codziennej pracy z~procesami.

\end{frame}
% ##################












% ######################################
\begin{frame}
  \frametitle{Uwagi techniczne}




\end{frame}
% ######################################










% ################################################
\section{Polecenie \texttt{grep} i~wyrażenia regularne}
% ################################################


% ######################################
\begin{frame}
  \frametitle{Polecenie \texttt{grep} i~wyrażenia regularne}




\end{frame}
% ######################################





% ######################################
\begin{frame}
  \frametitle{????}





\end{frame}
% ######################################





% ######################################
\begin{frame}
  \frametitle{}



\end{frame}
% ######################################




% ##################
\begin{frame}
  \frametitle{Problemy z~wyrażeniami regularnymi}




\end{frame}
% ##################





% ##################
\begin{frame}
  \frametitle{Mem dla ilustracji problemu}



\end{frame}
% ##################





% ##################
\begin{frame}
  \frametitle{Co ja robię?}



\end{frame}
% ##################





% ##################
\begin{frame}
  \frametitle{Praca z~\texttt{grep}em}




\end{frame}
% ##################





% ##################
\begin{frame}
  \frametitle{Praca z~\texttt{grep}em}



\end{frame}
% ##################





% ##################
\begin{frame}
  \frametitle{Praca z~\texttt{grep}em}




\end{frame}
% ##################











% ##################
\begin{frame}
  \frametitle{Wyszukiwanie całych wyrazów}




\end{frame}
% ##################





% ##################
\begin{frame}
  \frametitle{Wyszukiwanie całych wyrazów}




\end{frame}
% ##################
































% ######################################
\section{Polecenie \texttt{grep} i~wyrażenia regularne}
% ######################################


% ##################
\begin{frame}
  \frametitle{Wyrażenia regularne \texttt{grep}a}




\end{frame}
% ##################





% ##################
\begin{frame}
  \frametitle{Wyrażenia regularne \texttt{grep}a}


\end{frame}
% ##################





% ##################
\begin{frame}
  \frametitle{Wyrażenia regularne \texttt{grep}a}




\end{frame}
% ##################





% ##################
\begin{frame}
  \frametitle{Wyrażenia regularne \texttt{grep}a}



\end{frame}
% ##################





% ##################
\begin{frame}
  \frametitle{Funkcjonalności \texttt{grep}a można łączyć}



\end{frame}
% ##################





% ##################
\begin{frame}
  \frametitle{Rozbudowane wyrażenia regularne}



\end{frame}
% ##################







% % ##################
% \begin{frame}
%   \frametitle{}


% \end{frame}
% % ##################










% % % ##################
% % \begin{frame}
% %   \frametitle{}




% % \end{frame}
% % % ##################











% % % ##################
% % \begin{frame}
% %   \frametitle{}




% % \end{frame}
% % % ##################





% % % ##################
% % \begin{frame}
% %   \frametitle{}



% % \end{frame}
% % % ##################










% % % ######################################
% % \section{}
% % % ######################################


% % % ##################
% % \begin{frame}
% %   \frametitle{}



% % \end{frame}
% % % ##################







% % % ######################################
% % \section{}
% % % ######################################


% % % ##################
% % \begin{frame}
% %   \frametitle{}




% % \end{frame}
% % % ##################





% % % ##################
% % \begin{frame}
% %   \frametitle{}



% % \end{frame}
% % % ##################





% % % ##################
% % \begin{frame}
% %   \frametitle{}



% % \end{frame}
% % % ##################





% % % ##################
% % \begin{frame}
% %   \frametitle{}



% % \end{frame}
% % % ##################





% % % ##################
% % \begin{frame}
% %   \frametitle{}


% % \end{frame}
% % % ##################













% % % ##################
% % \begin{frame}
% %   \frametitle{}



% % \end{frame}
% % % ##################





% % % ##################
% % \begin{frame}
% %   \frametitle{}




% % \end{frame}
% % % ##################





% % % ##################
% % \begin{frame}
% %   \frametitle{}



% % \end{frame}
% % % ##################










% % % ######################################
% % \section{}
% % % ######################################


% % % ##################
% % \begin{frame}
% %   \frametitle{}



% % \end{frame}
% % % ##################





% % % ##################
% % \begin{frame}
% %   \frametitle{}



% % \end{frame}
% % % ##################







% % % ##################
% % \begin{frame}
% %   \frametitle{}



% % \end{frame}
% % % ##################





% % % ##################
% % \begin{frame}
% %   \frametitle{}




% % \end{frame}
% % % ##################





% % % ##################
% % \begin{frame}
% %   \frametitle{}



% % \end{frame}
% % % ##################





% % % ##################
% % \begin{frame}
% %   \frametitle{}



% % \end{frame}
% % % ##################





% % ##################
% \begin{frame}
%   \frametitle{????}




% \end{frame}
% % ##################





% % ##################
% \begin{frame}
%   \frametitle{????}




% \end{frame}
% % ##################





% % ##################
% \begin{frame}
%   \frametitle{????}




% \end{frame}
% % ##################





% % ##################
% \begin{frame}
%   \frametitle{????}




% \end{frame}
% % ##################






% % ##################
% \begin{frame}
%   \frametitle{?????}




% \end{frame}
% % ##################





% % ##################
% \begin{frame}
%   \frametitle{?????}




% \end{frame}
% % ##################





% % ##################
% \begin{frame}
%   \frametitle{?????}




% \end{frame}
% % ##################





% % ##################
% \begin{frame}
%   \frametitle{?????}



% \end{frame}
% % ##################





% % ##################
% \begin{frame}
%   \frametitle{?????}



% \end{frame}
% % ##################





% % ##################
% \begin{frame}
%   \frametitle{?????}



% \end{frame}
% % ##################





% % ##################
% \begin{frame}
%   \frametitle{????}



% \end{frame}
% % ##################





% % ##################
% \begin{frame}
%   \frametitle{?????}



% \end{frame}
% % ##################





% % ##################
% \begin{frame}
%   \frametitle{????}



% \end{frame}
% % ##################











% % ##################
% \begin{frame}
%   \frametitle{????}




% \end{frame}
% % ##################





% ################################################
\section{Materiały z~których korzystaliśmy}
% ################################################


% ######################################
\begin{frame}
  \frametitle{Materiały z których korzystaliśmy}


  Przygotowując tą prezentację korzystaliśmy z~następujących materiałów.

  Anish Singh Walia, Pankaj Kumar, Manikandan Kurup, \textit{Mastering Grep
    command in~Linux/Unix: A~Beginner's Tutorial},
  \parencite{WKK-Mastering-Grep-command-in-Linux-Unix-ETC-Ver-2026}.

\end{frame}
% ######################################










% ################################################
\appendix
% ################################################





% ######################################
\EndingSlide{Dziękuję! Pytania?}
% ######################################











% ####################################################################
% ####################################################################
% Bibliography

\printbibliography





% ####################################################################
% End of the document

\end{document}
